\documentclass[12pt,a4paper]{article}

% --- Metadata (per course) ---
\def\courseshortheader{THỐNG KÊ BAYES}
\def\coverbookline{NHÓM 5 --- XÁC SUẤT VÀ THỐNG KÊ}
\def\covermaintitle{Thống kê Bayes}
\def\coversubtitle{Giáo án lý thuyết --- Hậu nghiệm, MCMC, so sánh mô hình}

% Shared preamble for nhom_5 (requires \courseshortheader before \input)
\usepackage[utf8]{inputenc}
\usepackage[vietnamese]{babel}
\usepackage{amsmath,amssymb,amsthm}
\usepackage{geometry}
\usepackage{enumitem}
\usepackage[most]{tcolorbox}
\usepackage{hyperref}
\usepackage{fancyhdr}
\usepackage{graphicx}
\usepackage{booktabs}
\usepackage{tikz}
\usetikzlibrary{calc,arrows.meta,decorations.pathreplacing,positioning}
\usepackage{eso-pic}
\usepackage{xcolor}

\geometry{a4paper, margin=2cm, bottom=2.5cm}
\hypersetup{colorlinks=true, linkcolor=blue!70!black, urlcolor=blue}
\setlist[itemize]{topsep=4pt, itemsep=2pt}
\setlist[enumerate]{topsep=4pt, itemsep=2pt}

\AddToShipoutPictureFG{%
  \AtPageCenter{%
    \begin{tikzpicture}[remember picture, overlay]
      \node[rotate=45, scale=1, opacity=0.06]{%
        \fontsize{60}{72}\selectfont\bfseries\sffamily Đặng Tấn Quí%
      };
    \end{tikzpicture}%
  }%
}

\pagestyle{fancy}
\fancyhf{}
\fancyhead[L]{\small\textit{\courseshortheader}}
\fancyhead[R]{\small\textit{Đặng Tấn Quí}}
\fancyfoot[C]{\thepage}
\fancyfoot[L]{\footnotesize\textcopyright\ Đặng Tấn Quí}
\fancyfoot[R]{\footnotesize SĐT: 0377\,122\,210}
\renewcommand{\headrulewidth}{0.4pt}
\renewcommand{\footrulewidth}{0.4pt}

\fancypagestyle{plain}{%
  \fancyhf{}
  \fancyfoot[C]{\thepage}
  \fancyfoot[L]{\footnotesize\textcopyright\ Đặng Tấn Quí}
  \fancyfoot[R]{\footnotesize SĐT: 0377\,122\,210}
  \renewcommand{\headrulewidth}{0pt}
  \renewcommand{\footrulewidth}{0.4pt}
}

\newtcolorbox{dongco}[1][]{%
  enhanced, breakable,
  colback=teal!4!white, colframe=teal!60!black,
  fonttitle=\bfseries, title={Động cơ học -- Tại sao cần khái niệm này?}, #1%
}
\newtcolorbox{ynghia}[1][]{%
  enhanced, breakable,
  colback=purple!3!white, colframe=purple!50!black,
  fonttitle=\bfseries, title={Ý nghĩa \& Trực giác}, #1%
}
\newtcolorbox{ungdung}[1][]{%
  enhanced, breakable,
  colback=olive!5!white, colframe=olive!60!black,
  fonttitle=\bfseries, title={Ứng dụng thực tế}, #1%
}
\newtcolorbox{dinhnghia}[1][]{%
  enhanced, breakable,
  colback=blue!3!white, colframe=blue!60!black,
  fonttitle=\bfseries, title={Định nghĩa}, #1%
}
\newtcolorbox{dinhly}[1][]{%
  enhanced, breakable,
  colback=green!3!white, colframe=green!50!black,
  fonttitle=\bfseries, title={Định lý}, #1%
}
\newtcolorbox{tinhchat}[1][]{%
  enhanced, breakable,
  colback=yellow!5!white, colframe=orange!50!black,
  fonttitle=\bfseries, title={Tính chất}, #1%
}
\newtcolorbox{phuongphap}[1][]{%
  enhanced, breakable,
  colback=gray!5!white, colframe=gray!50!black,
  fonttitle=\bfseries, title={Phương pháp}, #1%
}
\newtcolorbox{luuy}[1][]{%
  enhanced, breakable,
  colback=red!3!white, colframe=red!50!black,
  fonttitle=\bfseries, title={Lưu ý}, #1%
}
\newtcolorbox{congthuc}[1][]{%
  enhanced, breakable,
  colback=cyan!3!white, colframe=cyan!50!black,
  fonttitle=\bfseries, title={Công thức}, #1%
}
\newtcolorbox{vidu}[1][]{%
  enhanced, breakable,
  colback=violet!3!white, colframe=violet!50!black,
  fonttitle=\bfseries, title={Ví dụ}, #1%
}
\newtcolorbox{phanvidu}[1][]{%
  enhanced, breakable,
  colback=red!2!white, colframe=red!40!black,
  fonttitle=\bfseries, title={Phản ví dụ}, #1%
}
\newtcolorbox{tongket}[1][]{%
  enhanced, breakable,
  colback=blue!4!white, colframe=blue!70!black,
  fonttitle=\bfseries, title={Tổng kết chương}, #1%
}


\begin{document}

\thispagestyle{empty}
\begin{tikzpicture}[remember picture, overlay]
  \fill[blue!80!black] (current page.south west) rectangle (current page.north east);
  \fill[blue!60!cyan, opacity=0.8]
    (current page.south west) -- (current page.north west) --
    ([xshift=-3cm]current page.north east) -- (current page.south east) -- cycle;
  \foreach \r/\o in {6/0.05, 4.5/0.08, 3/0.12} {
    \fill[white, opacity=\o] ([xshift=2cm, yshift=-2cm]current page.north east) circle (\r cm);
  }
  \draw[white, line width=2pt] ([yshift=-5cm]current page.north west) ++(2cm,0) -- ++(17cm,0);
  \draw[white, line width=2pt] ([yshift=-16cm]current page.north west) ++(2cm,0) -- ++(17cm,0);
  \node[anchor=west, white, font=\fontsize{28}{34}\bfseries\selectfont]
    at ([xshift=3cm, yshift=-7.5cm]current page.north west) {\coverbookline};
  \node[anchor=west, white, font=\fontsize{20}{26}\selectfont]
    at ([xshift=3cm, yshift=-9cm]current page.north west) {\covermaintitle};
  \node[anchor=west, white, font=\fontsize{16}{22}\selectfont]
    at ([xshift=3cm, yshift=-10.2cm]current page.north west) {\coversubtitle};
  \node[anchor=west, white, font=\Large\bfseries]
    at ([xshift=3cm, yshift=-13cm]current page.north west) {Biên soạn:};
  \node[anchor=west, yellow!80!white, font=\fontsize{24}{30}\bfseries\selectfont]
    at ([xshift=3cm, yshift=-14.5cm]current page.north west) {ĐẶNG TẤN QUÍ};
  \node[anchor=west, white!90, font=\large]
    at ([xshift=3cm, yshift=-18cm]current page.north west) {SĐT: 0377\,122\,210};
  \node[anchor=south, white!80, font=\small]
    at ([yshift=1.5cm]current page.south) {Tài liệu có bản quyền --- Cấm sao chép khi chưa được sự đồng ý của tác giả};
  \node[anchor=south east, white, font=\Large\bfseries]
    at ([xshift=-2cm, yshift=3cm]current page.south east) {\the\year};
\end{tikzpicture}

\newpage
\tableofcontents
\newpage

\section*{Lời nói đầu}
\addcontentsline{toc}{section}{Lời nói đầu}

Tài liệu thuộc \textbf{Nhóm 5 --- Xác suất và thống kê}, biên soạn theo cùng triết lý với giáo án Đại số tuyến tính: học để \emph{hiểu} và \emph{dùng}, không chỉ để ghi nhớ công thức.

\begin{enumerate}
  \item \textbf{Động cơ trước định nghĩa.} Mỗi phần lớn mở bằng ô \textsf{\color{teal!60!black} Động cơ học} --- câu hỏi thực tế hoặc bài toán mẫu cần khái niệm đó.
  \item \textbf{Trực giác trước công thức.} Ô \textsf{\color{purple!50!black} Ý nghĩa \& Trực giác} giải thích ``công thức đang nói gì'' (xác suất = tần suất dài hạn, kỳ vọng = trung bình có trọng số, v.v.).
  \item \textbf{Ví dụ có kiểm tra.} Ô \textsf{\color{violet!50!black} Ví dụ} nên tự làm trước khi đọc lời giải; phần thống kê ứng dụng cần gắn dữ liệu số cụ thể.
  \item \textbf{Tổng kết cuối mỗi chương.} Ô \textsf{\color{blue!70!black} Tổng kết chương} liệt kê kết quả phải nhớ và liên kết sang phần sau.
  \item \textbf{Phân tầng theo môn.} X1--X3 là nền; X4--X5 nâng cao lý thuyết; X6--X8 chuyên sâu (đa biến, Bayes, quá trình ngẫu nhiên).
\end{enumerate}

\noindent\textbf{Cách đọc:} Động cơ $\to$ Định nghĩa / Định lý $\to$ Ví dụ $\to$ Ý nghĩa $\to$ Tổng kết. Với môn thống kê, luôn hỏi: \emph{mẫu là gì, tham số nào chưa biết, giả thuyết $H_0$ là gì}.

\newpage


\section{Mở đầu}

\begin{dongco}[title={Động cơ --- Mở đầu}]
Thống kê cổ điển coi $\theta$ cố định; thống kê Bayes coi $\theta$ ngẫu nhiên với phân phối tiên nghiệm, cập nhật bằng dữ liệu thành \emph{hậu nghiệm}.
\end{dongco}

\newpage

\section{Nguyên lý Bayes}

\begin{dongco}[title={Động cơ --- Nguyên lý Bayes}]
Công thức $p(	heta|data) \propto p(data|	heta)p(	heta)$ là trục của toàn bộ môn.
\end{dongco}

\begin{dinhnghia}[title={Suy luận Bayes}]
  Cho dữ liệu $\mathbf{x}$ và tham số $\theta$:
  \[
    \underbrace{p(\theta|\mathbf{x})}_{\text{posterior}} = \frac{\overbrace{p(\mathbf{x}|\theta)}^{\text{likelihood}}\;\overbrace{p(\theta)}^{\text{prior}}}{\underbrace{p(\mathbf{x})}_{\text{marginal}}} \propto p(\mathbf{x}|\theta)\,p(\theta).
  \]
\end{dinhnghia}

\begin{dinhnghia}[title={Prior (phân phối tiên nghiệm)}]
  \begin{itemize}
    \item \textbf{Thông tin} (informative): phản ánh kiến thức trước.
    \item \textbf{Không thông tin} (non-informative): ít ảnh hưởng, ví dụ: đều, Jeffreys.
    \item \textbf{Prior Jeffreys:} $p(\theta) \propto \sqrt{I(\theta)}$ (bất biến tham số hóa).
  \end{itemize}
\end{dinhnghia}

\begin{dinhnghia}[title={Ước lượng Bayes}]
  \begin{itemize}
    \item \textbf{MAP} (Maximum A Posteriori): $\hat{\theta}_{\text{MAP}} = \arg\max_\theta p(\theta|\mathbf{x})$.
    \item \textbf{Trung bình hậu nghiệm}: $\hat{\theta}_{\text{Bayes}} = E[\theta|\mathbf{x}]$ (cực tiểu rủi ro bình phương).
    \item \textbf{Trung vị hậu nghiệm}: cực tiểu rủi ro trị tuyệt đối.
  \end{itemize}
\end{dinhnghia}

\begin{dinhnghia}[title={Khoảng tin cậy Bayes (credible interval)}]
  $C$ sao cho $P(\theta \in C|\mathbf{x}) = 1 - \alpha$. \textbf{HPD} (Highest Posterior Density): khoảng ngắn nhất.
\end{dinhnghia}

%% ================================================================
%%  CHƯƠNG 2: HỌ PHÂN PHỐI LIÊN HỢP
%% ================================================================

\begin{tongket}[title={Tổng kết: Nguyên lý Bayes}]
\begin{enumerate}
  \item Nắm lại các \textbf{định nghĩa} và \textbf{công thức} cốt lõi trong mục ``Nguyên lý Bayes''.
  \item Tự làm lại ít nhất một ví dụ (nếu có) hoặc áp dụng công thức vào một tình huống số.
  \item Ghi chú điều kiện áp dụng (giả thiết độc lập, phân phối chuẩn, $n$ đủ lớn, v.v.).
\end{enumerate}
\end{tongket}

\section{Họ phân phối liên hợp}

\begin{dongco}[title={Động cơ --- Họ phân phối liên hợp}]
Chọn prior--likelihood sao cho hậu nghiệm thuộc cùng họ --- tính được đóng.
\end{dongco}

\begin{dinhnghia}
  Họ prior $\mathcal{P}$ là \textbf{liên hợp} (conjugate) với likelihood nếu posterior cũng thuộc $\mathcal{P}$.
\end{dinhnghia}

\begin{congthuc}[title={Các cặp liên hợp quan trọng}]
  \begin{itemize}
    \item \textbf{Bernoulli/Nhị thức + Beta}: $\text{Beta}(\alpha, \beta) \to \text{Beta}(\alpha + k, \beta + n - k)$.
    \item \textbf{Poisson + Gamma}: $\text{Gamma}(\alpha, \beta) \to \text{Gamma}(\alpha + \sum x_i, \beta + n)$.
    \item \textbf{Chuẩn (biết $\sigma^2$) + Chuẩn}: $\mathcal{N}(\mu_0, \tau_0^2) \to \mathcal{N}(\mu_n, \tau_n^2)$, với $\tau_n^{-2} = \tau_0^{-2} + n/\sigma^2$.
    \item \textbf{Chuẩn (biết $\mu$) + Inverse-Gamma}: cho $\sigma^2$.
    \item \textbf{Đa thức + Dirichlet}: $\text{Dir}(\boldsymbol{\alpha}) \to \text{Dir}(\boldsymbol{\alpha} + \mathbf{n})$.
  \end{itemize}
\end{congthuc}

\begin{luuy}
  Posterior liên hợp = ``trung bình có trọng số'' giữa prior và dữ liệu. Khi $n \to \infty$, posterior tập trung quanh MLE (posterior consistency).
\end{luuy}

%% ================================================================
%%  CHƯƠNG 3: PHƯƠNG PHÁP MCMC
%% ================================================================

\begin{tongket}[title={Tổng kết: Họ phân phối liên hợp}]
\begin{enumerate}
  \item Nắm lại các \textbf{định nghĩa} và \textbf{công thức} cốt lõi trong mục ``Họ phân phối liên hợp''.
  \item Tự làm lại ít nhất một ví dụ (nếu có) hoặc áp dụng công thức vào một tình huống số.
  \item Ghi chú điều kiện áp dụng (giả thiết độc lập, phân phối chuẩn, $n$ đủ lớn, v.v.).
\end{enumerate}
\end{tongket}

\section{Phương pháp MCMC}

\begin{dongco}[title={Động cơ --- Phương pháp MCMC}]
Khi hậu nghiệm không tính được, lấy mẫu xấp xỉ bằng chuỗi Markov.
\end{dongco}

\begin{dinhnghia}[title={Monte Carlo Markov Chain}]
  Xây dựng xích Markov $(\theta^{(1)}, \theta^{(2)}, \ldots)$ có phân phối dừng là $p(\theta|\mathbf{x})$. Sau ``burn-in'', mẫu xấp xỉ từ posterior.
\end{dinhnghia}

\begin{phuongphap}[title={Metropolis--Hastings}]
  \begin{enumerate}
    \item Đề xuất $\theta^* \sim q(\cdot|\theta^{(t)})$.
    \item Tính xác suất chấp nhận: $\alpha = \min\!\left(1,\; \dfrac{p(\theta^*|\mathbf{x})\,q(\theta^{(t)}|\theta^*)}{p(\theta^{(t)}|\mathbf{x})\,q(\theta^*|\theta^{(t)})}\right)$.
    \item Chấp nhận $\theta^{(t+1)} = \theta^*$ với xác suất $\alpha$; nếu không, $\theta^{(t+1)} = \theta^{(t)}$.
  \end{enumerate}
\end{phuongphap}

\begin{phuongphap}[title={Gibbs Sampling}]
  Trường hợp đặc biệt M--H: lần lượt lấy mẫu từng thành phần theo phân phối có điều kiện đầy đủ:
  \[
    \theta_j^{(t+1)} \sim p(\theta_j|\theta_1^{(t+1)}, \ldots, \theta_{j-1}^{(t+1)}, \theta_{j+1}^{(t)}, \ldots, \theta_p^{(t)}, \mathbf{x}).
  \]
  Xác suất chấp nhận luôn bằng 1.
\end{phuongphap}

\begin{phuongphap}[title={Chẩn đoán MCMC}]
  \begin{itemize}
    \item Trace plot: kiểm tra hội tụ trực quan.
    \item $\hat{R}$ (Gelman--Rubin): so sánh phương sai giữa/trong xích, $\hat{R} \approx 1$ tốt.
    \item Effective sample size (ESS): $\text{ESS} = n / (1 + 2\sum \rho_k)$.
    \item Autocorrelation: thinning nếu tự tương quan cao.
  \end{itemize}
\end{phuongphap}

%% ================================================================
%%  CHƯƠNG 4: SO SÁNH MÔ HÌNH BAYES
%% ================================================================

\begin{tongket}[title={Tổng kết: Phương pháp MCMC}]
\begin{enumerate}
  \item Nắm lại các \textbf{định nghĩa} và \textbf{công thức} cốt lõi trong mục ``Phương pháp MCMC''.
  \item Tự làm lại ít nhất một ví dụ (nếu có) hoặc áp dụng công thức vào một tình huống số.
  \item Ghi chú điều kiện áp dụng (giả thiết độc lập, phân phối chuẩn, $n$ đủ lớn, v.v.).
\end{enumerate}
\end{tongket}

\section{So sánh mô hình Bayes}

\begin{dongco}[title={Động cơ --- So sánh mô hình Bayes}]
Bayes factor, DIC, WAIC --- chọn mô hình có dự đoán tốt hơn.
\end{dongco}

\begin{dinhnghia}[title={Bayes factor}]
  So sánh $M_1$ vs $M_2$:
  \[
    BF_{12} = \frac{p(\mathbf{x}|M_1)}{p(\mathbf{x}|M_2)}, \qquad p(\mathbf{x}|M_k) = \int p(\mathbf{x}|\theta, M_k)\,p(\theta|M_k)\,d\theta.
  \]
  $BF_{12} > 1$: dữ liệu ủng hộ $M_1$. $BF > 10$: ủng hộ mạnh.
\end{dinhnghia}

\begin{dinhnghia}[title={Tiêu chuẩn thông tin}]
  \begin{itemize}
    \item \textbf{BIC:} $-2\ln L(\hat{\theta}) + k\ln n$ (xấp xỉ $-2\ln BF$).
    \item \textbf{DIC:} $D(\bar{\theta}) + 2p_D$ (số tham số hiệu quả $p_D$).
    \item \textbf{WAIC:} cải tiến DIC, dùng posterior pointwise.
  \end{itemize}
\end{dinhnghia}

\begin{phuongphap}[title={Dự báo Bayes}]
  \textbf{Phân phối dự báo hậu nghiệm} (posterior predictive):
  \[
    p(\tilde{x}|\mathbf{x}) = \int p(\tilde{x}|\theta)\,p(\theta|\mathbf{x})\,d\theta.
  \]
  Tích hợp bất định tham số; khoảng dự báo tự nhiên.
\end{phuongphap}

\begin{tongket}[title={Tổng kết: So sánh mô hình Bayes}]
\begin{enumerate}
  \item Nắm lại các \textbf{định nghĩa} và \textbf{công thức} cốt lõi trong mục ``So sánh mô hình Bayes''.
  \item Tự làm lại ít nhất một ví dụ (nếu có) hoặc áp dụng công thức vào một tình huống số.
  \item Ghi chú điều kiện áp dụng (giả thiết độc lập, phân phối chuẩn, $n$ đủ lớn, v.v.).
\end{enumerate}
\end{tongket}



\end{document}
